\chapter{Background and Related Work}
\label{ch:background}

\epigraph{There is no AI without PI.}{\textit{Wil van der Aalst}, on
object-centric process mining as the enabler of artificial intelligence
\citep{noai-nopi-2025}}

This chapter establishes the four foundations on which the thesis rests: process
mining and the paradigm founded by van der Aalst (\cref{sec:bg-pm}); the Language
Server Protocol (\cref{sec:bg-lsp}); the Model Context Protocol
(\cref{sec:bg-mcp}); and the Agent-to-Agent protocol together with the wider
agent-interoperability landscape (\cref{sec:bg-a2a,sec:bg-landscape}).
\Cref{sec:bg-synthesis} draws them together into the observation that organises
the rest of the thesis: each is the same decoupling, performed at a different
layer.

\section{Process Mining and the van der Aalst Paradigm}
\label{sec:bg-pm}

\subsection{Origins and definition}

Process mining sits between data mining and business process management. Its
canonical statement, from the \emph{Process Mining Manifesto} issued by the IEEE
Task Force on Process Mining---a document drafted by more than seventy-five people
across more than fifty organisations \citep{tfpm-manifesto}---is that ``the basic
idea of process mining is to discover, monitor and improve real processes
(i.e., not assumed processes) by extracting knowledge from event logs readily
available in today's (information) systems'' \citep{pm-manifesto-2011}. The
field's reference monograph is van der Aalst's \emph{Process Mining: Data Science
in Action} \citep{vdaalst-book-2016}.

The discipline was, in a real sense, founded and sustained by a single central
figure. Van der Aalst holds the Chair of Process and Data Science (PADS) at RWTH
Aachen University, is part-time affiliated with Fraunhofer FIT, and since August
2021 has been Chief Scientist at Celonis \citep{vdaalst-cv,celonis-godfather}; he
is a Fellow of the IFIP, IEEE, and ACM, and received the Alexander von Humboldt
Professorship in 2018 \citep{vdaalst-cv}. He is widely styled ``the Godfather of
Process Mining''---a descriptor used both on his own curriculum vitae and in the
official Celonis appointment \citep{vdaalst-cv,celonis-godfather}---and is among
the most highly cited computer scientists, with an $h$-index of roughly
$188$--$191$ and on the order of $170{,}000$ citations as reported by Google
Scholar in the mid-2020s.\footnote{Citation metrics are time-variant; the figures
are reported as a snapshot rather than a fixed quantity \citep{vdaalst-cv}.} We
foreground this not as flattery but because the thesis argues that his
paradigm---not merely his tools---generalises beyond its origin domain, and the
strength of that claim is proportional to the maturity of what is being
generalised.

\subsection{The three types of process mining}

Process mining comprises three complementary activities
\citep{pm-manifesto-2011,vdaalst-book-2016}:

\begin{description}[leftmargin=1.4em,style=nextline]
  \item[Discovery] taking an event log and producing a process model without any
  a-priori model.
  \item[Conformance checking] comparing an existing model with an event log of
  the same process to locate where reality and model agree and where they
  diverge.
  \item[Enhancement] extending or improving a model using information recorded in
  the actual log.
\end{description}

Conformance checking is the activity this thesis elevates to a protocol surface.

\subsection{Conformance checking and the four quality dimensions}
\label{sec:bg-conformance}

Two principal techniques realise conformance checking. \emph{Token-based replay}
replays each trace against a (Petri-net) model, counting produced, consumed,
missing, and remaining tokens to compute a fitness score; \emph{alignments}, the
more modern technique, compute an optimal correspondence between trace and model,
yielding more precise diagnostics at higher computational cost
\citep{conformance-review-2020}. The quality of a discovered or checked model is
judged along four dimensions that are in tension with one another
\citep{vdaalst-book-2016}:

\begin{description}[leftmargin=1.4em,style=nextline]
  \item[Fitness] the model can replay the observed behaviour.
  \item[Precision] the model does not allow much behaviour that was never
  observed (avoiding underfitting).
  \item[Generalisation] the model is not over-specialised to the log (avoiding
  overfitting).
  \item[Simplicity] the model is as simple as the behaviour permits (Occam's
  razor).
\end{description}

These four dimensions recur in \cref{ch:framework} as the criteria a conformance
surface for \emph{any} language must balance.

\subsection{Event-log standards: XES, OCEL, and OCEL 2.0}
\label{sec:bg-ocel}

Conformance requires logs, and logs require a standard. The first was XES
(eXtensible Event Stream), ratified as IEEE Std~1849-2016 \citep{xes-2016}. XES
assumes a single \emph{case notion}: every event belongs to exactly one case.
Real systems violate this assumption---an event may touch an order, several
items, a package, and a customer at once---and forcing a single case identifier
produces two pathologies that van der Aalst names \emph{convergence} (an event is
duplicated across cases) and \emph{divergence} (ordering within a case is lost)
\citep{ocpn-2020}.

The response is the \emph{Object-Centric Event Log} (OCEL): OCEL~1.0 appeared in
2020, and OCEL~2.0 was specified in 2023--2024 \citep{ocel2-2024,ocel-standard}.
Per its specification, OCEL~2.0 ``forms the new, more expressive standard'' and,
unlike its predecessor, ``can depict changes in objects, provide information on
object relationships, and qualify these relationships'' \citep{ocel2-2024}.
Concretely it adds dynamic object attributes (values that change over time),
object-to-object relationships (two objects related without sharing an event),
and qualified relationships (the role an object plays in an event), with exchange
formats in SQLite, XML, and JSON \citep{ocel2-2024,ocel-standard}. OCEL~2.0 is
the data model the artifact of this thesis emits, and \cref{ch:all-language}
argues it is the right abstraction for \emph{all} language, not only business
process.

\subsection{Object-centric process mining and ``no AI without PI''}
\label{sec:bg-ocpm}

Object-centric process mining (OCPM) generalises discovery and conformance to
logs with many interrelated object types; object-centric Petri nets give each
object type its own subnet within one integrated net \citep{ocpn-2020}. Van der
Aalst frames OCPM as ``the next frontier'' and, most recently, as the
indispensable substrate for artificial intelligence: in \emph{No AI Without PI!}
he argues that object-centric process mining is ``the missing link connecting
data and processes'' and introduces \emph{Process Intelligence} (PI) as the
amalgamation of process-centric, data-driven techniques required to ground
generative, predictive, and prescriptive AI \citep{noai-nopi-2025}. A parallel
strand in the community has begun to re-examine process mining in light of
LLM-based agents, proposing agent-workflow decompositions of mining tasks
\citep{rethinking-pm-agents-2024}. This thesis is, in effect, the contrapositive
of van der Aalst's slogan applied to language: \emph{no conformance for all
language without an object-centric event abstraction}.

\subsection{Methodology: PM\texorpdfstring{$^2$}{2} and L*}

Process-mining projects are themselves structured. The PM\textsuperscript{2}
methodology defines staged activities with explicit inputs and outputs and was
validated on an industrial purchasing process \citep{pm2-2015}; the earlier L*
life-cycle model in \citet{vdaalst-book-2016} sets out a five-stage progression
from planning and extraction through model construction to operational support.
\Cref{ch:methodology} positions the present work's design-science process
alongside these.

\section{The Language Server Protocol}
\label{sec:bg-lsp}

LSP was announced on 27~June~2016 as a collaboration of Red Hat, Codenvy, and
Microsoft, having grown out of the integration of the OmniSharp (C\#) and
TypeScript servers into VS~Code \citep{lsp-redhat-2016}. Its design is built on
JSON-RPC~2.0: the content of every message ``uses JSON-RPC 2.0 to describe
requests, responses and notifications'', framed by a required \code{Content-Length}
header \citep{lsp-spec-318}. Communication begins with capability negotiation:
the \code{initialize} request, ``the first request \dots\ from the client to the
server'', exchanges client and server capabilities, and features may further be
registered dynamically thereafter \citep{lsp-spec-318}.

The protocol's load-bearing idea is captured in \cref{sec:motivation}: it
decouples language intelligence from the editor so that the integration burden
falls from $M \times N$ to $M + N$ \citep{lsp-official,lsp-mxn-fcc}. Versions
3.16 (December~2020) and 3.17 (May~2022) added semantic tokens, call and type
hierarchies, inlay hints, and notebook support; version 3.18 is, at the time of
writing, explicitly ``under development'' rather than finalised, a status this
thesis is careful to preserve \citep{lsp-spec-317,lsp-spec-318}.

Crucially for \cref{ch:all-language}, LSP has already escaped the boundary of
programming languages in practice. LTeX is a language server whose ``language''
is natural language: its documentation states that ``the `language' is not a
single natural language, but all natural languages supported by LanguageTool''
\citep{ltex}. Other servers target prose (textLSP) and data/configuration formats
(the Red Hat YAML language server) \citep{textlsp,yamlls}. These are existence
proofs, not a theory; supplying the theory---a criterion for what counts as a
servable language---is contribution~C2.

\section{The Model Context Protocol}
\label{sec:bg-mcp}

MCP was open-sourced by Anthropic on 25~November~2024 as ``a new standard for
connecting AI assistants to the systems where data lives''
\citep{mcp-anthropic-2024}. Like LSP, it is a JSON-RPC~2.0 protocol with explicit
capability negotiation \citep{mcp-spec-2025}, and it defines three roles: hosts
(``LLM applications that initiate connections''), clients (``connectors within
the host application''), and servers (``services that provide context and
capabilities'') \citep{mcp-spec-2025,mcp-architecture}. Servers expose three
primitives---\emph{resources} (context and data), \emph{prompts} (templated
workflows), and \emph{tools} (functions the model may execute)---while clients may
offer \emph{sampling}, \emph{roots}, and \emph{elicitation} back to servers
\citep{mcp-spec-2025}. Its transports are stdio for local processes and, since
the 2025-03-26 revision, ``Streamable HTTP'', which replaced the earlier HTTP+SSE
transport \citep{mcp-transports-2025,mcp-architecture}.

MCP's adoption through 2025 was unusually rapid for a vendor-originated standard:
OpenAI announced support on 26~March~2025 \citep{openai-mcp-techcrunch}, Google
followed on 9~April~2025 \citep{google-mcp-techcrunch}, and Microsoft integrated
it into Copilot Studio and Windows. In December~2025 Anthropic donated MCP to the
Agentic AI Foundation, ``a directed fund under the Linux Foundation'', co-founded
with Block and OpenAI and backed by AWS, Bloomberg, Cloudflare, Google, and
Microsoft, while stating that the project's governance model would remain
unchanged \citep{mcp-aaif-2025,lf-aaif-press}.

\section{The Agent-to-Agent Protocol}
\label{sec:bg-a2a}

Google published A2A on 9~April~2025 with more than fifty technology partners
\citep{a2a-google-2025}. Its central artefact is the \emph{Agent Card}, a JSON
document advertising an agent's identity, capabilities, skills, endpoint, and
authentication requirements, served---in current revisions---at the well-known
path \code{/.well-known/agent-card.json} \citep{a2a-spec,a2a-discovery}. The unit
of work is the \emph{Task}, with a defined lifecycle (\code{submitted},
\code{working}, \code{input-required}, \code{completed}, \code{failed},
\code{canceled}, and related states); communication uses \emph{Messages} and
\emph{Artifacts} composed of typed \emph{Parts} \citep{a2a-spec}. Like LSP and
MCP, A2A reuses ``existing, well-understood standards (HTTP, JSON-RPC 2.0,
Server-Sent Events)'' \citep{a2a-spec}.

A2A is positioned as complementary to, not competitive with, MCP. The
specification states that ``the Model Context Protocol (MCP) defines how an AI
agent interacts with \dots\ tools and resources'' while ``the Agent2Agent
Protocol focuses on enabling different agents to collaborate'': an application
uses A2A to talk to other agents, each of which uses MCP internally for its tools
\citep{a2a-mcp}. In June~2025 Google donated A2A to the Linux Foundation, seeding
an Agent2Agent project with AWS, Cisco, Google, Microsoft, Salesforce, SAP, and
ServiceNow \citep{a2a-lf-donation,lf-a2a-press}.

\section{The Agent-Interoperability Landscape}
\label{sec:bg-landscape}

A2A and MCP are the most prominent but not the only agent-interoperability
efforts. IBM's Agent Communication Protocol (ACP), a REST-native alternative,
merged with A2A under the Linux Foundation in August~2025 \citep{acp-a2a-merge};
the Agent Network Protocol (ANP) pursues a decentralised ``HTTP of the agentic
web'' built on decentralised identifiers \citep{anp-github}; and Cisco's AGNTCY
collective promotes an ``Internet of Agents'' spanning discovery, composition,
and evaluation \citep{agntcy-cisco}. Commentary oscillates between a ``protocol
wars'' framing and a ``co-opetition'' one, but the institutional trajectory---A2A
and then MCP arriving under the Linux Foundation umbrella within six months
\citep{lf-a2a-press,lf-aaif-press}---favours consolidation. The emerging
reference architecture is a two-layer stack: MCP for vertical agent-to-tool
integration, A2A for horizontal agent-to-agent coordination
\citep{a2a-google-2025,a2a-mcp}. This thesis adds a third, lower stratum---LSP,
generalised---and a cross-cutting surface---conformance.

\section{Synthesis: Three Decouplings, One Trajectory}
\label{sec:bg-synthesis}

\Cref{tab:three-decouplings} reads the three protocols as instances of one move.
Each isolates a capability that was previously fused to a host, exposes it behind
a JSON-RPC contract with capability negotiation, and thereby converts an
$M \times N$ integration burden into $M + N$.

\begin{table}[t]
\centering
\caption[Three decouplings, one trajectory]{LSP, MCP, and A2A as three instances
of the same decoupling. Each exposes a previously host-bound capability behind a
JSON-RPC contract with capability negotiation, collapsing $M \times N$ into
$M + N$.}
\label{tab:three-decouplings}
\small
\begin{tabularx}{\textwidth}{@{}l l l Y@{}}
\toprule
\textbf{Protocol} & \textbf{Year} & \textbf{Decouples\,\dots} & \textbf{\dots from} \\
\midrule
LSP & 2016 & language intelligence & the editor \citep{lsp-redhat-2016} \\
MCP & 2024 & context and tools & the model \citep{mcp-anthropic-2024} \\
A2A & 2025 & agent capability & the framework \citep{a2a-google-2025} \\
\bottomrule
\end{tabularx}
\end{table}

Three facts make the ``one trajectory'' reading more than analogy. First, common
mechanism: all three are JSON-RPC protocols with explicit client/server
capability negotiation \citep{lsp-spec-318,mcp-spec-2025,a2a-spec}. Second,
acknowledged lineage: the MCP specification records that it ``takes some
inspiration from the Language Server Protocol'' \citep{mcp-spec-2025}, and A2A
defines itself relative to MCP \citep{a2a-mcp}. Third, converging governance: MCP
and A2A now share a foundation \citep{lf-a2a-press,lf-aaif-press}. What the stack
still lacks---and what process mining supplies---is a surface that asks not
whether a message was delivered but whether the resulting behaviour conformed.
\Cref{ch:framework} builds that surface.
